% !TEX root = ../main.tex
\subsection{Setup}
We evaluate on a discrete-event simulator of a single host with $N{=}16$
pinned worker cores. New requests are randomly assigned to cores; only
waiting requests migrate; running requests do not. Unless noted, the
defaults are $check{=}1\,\mu s$, $K{=}16$, $H{=}4$, $B{=}1$,
$\epsilon{=}2\,\mu s$, $\theta{=}0$, and $M{=}0.5\,\mu s$. Workloads are
W1 (near-saturation), W2 (MMPP burst-skew), and W3 (heavy-tailed).
Baselines are RSS-local FIFO (d-FCFS), work stealing, intra-host
proactive migration, and threshold/shortest-queue migration. All numbers
are medians over 5 seeds unless a 10-seed CI is reported.

\note{Regenerate Table~\ref{tab:main} and figures from
artifacts/step-15-rescuesched and step-17-rescuesched-closure via
scripts/rescue\_analysis.py before submission.}

\begin{table}[t]
\centering
\caption{W3 heavy-tail, median over 5 seeds. RescueSched lowers SLO
violation while migrating fewer requests than proactive migration.}
\label{tab:main}
\small
\begin{tabular}{@{}llrrrr@{}}
\toprule
$\rho$ & Method & SLO viol. & Move rate & BMR & UMR \\
\midrule
0.85 & Work stealing        & 0.203 & 0.328 & 0.00 & 0.00 \\
0.85 & Proactive mig.\ (M0) & 0.154 & 0.400 & 0.00 & 0.00 \\
0.85 & \textbf{RescueSched} & \textbf{0.094} & 0.286 & 1.00 & 0.00 \\
\bottomrule
\end{tabular}
\end{table}

\subsection{Main Result}
On W3 at $\rho{=}0.85$ RescueSched cuts the SLO violation ratio to
$0.094$ from $0.203$ (work stealing) and $0.154$ (proactive migration),
while its move rate ($0.286$) is \emph{lower} than proactive migration's
($0.400$)---the gain is quality, not volume. Its P99/P99.9 are not always
the lowest; RescueSched optimizes SLO rescue and migration usefulness,
not raw tail latency, and we state this trade-off explicitly.

\subsection{Ablations}
Removing rescuability (\emph{NoRescuable}) migrates far more
($0.60$ vs $0.29$) but collapses BMR to $0.31$, raises UMR to $0.69$, and
degrades SLO violation to $0.169$---confirming that task-level
rescuability is necessary. Removing target safety leaves \emph{actual}
harmful migrations at $0$ under the current FIFO append-to-tail model; it
reduces only \emph{predicted} target-risk exposure. We report this as a
model boundary rather than a strong claim.

\subsection{Sensitivity and Boundaries}
$check{=}1\,\mu s$ is the operating point; $5\,\mu s$ misses the rescue
window. Budget $B{=}1$ dominates $B{=}2,4$, which add unsaved migrations.
Under W1 global overload ($\rho{=}0.95$) RescueSched performs $0$
migrations, confirming it does not fabricate capacity.